
\documentclass{article}
\usepackage{colortbl}
\usepackage{makecell}
\usepackage{multirow}
\usepackage{supertabular}

\begin{document}

\newcounter{utterance}

\twocolumn

{ \footnotesize  \setcounter{utterance}{1}
\setlength{\tabcolsep}{0pt}
\begin{supertabular}{c@{$\;$}|p{.15\linewidth}@{}p{.15\linewidth}p{.15\linewidth}p{.15\linewidth}p{.15\linewidth}p{.15\linewidth}}

    \# & $\;$A & \multicolumn{4}{c}{Game Master} & $\;\:$B\\
    \hline 

    \theutterance \stepcounter{utterance}  

    & & \multicolumn{4}{p{0.6\linewidth}}{\cellcolor[rgb]{0.9,0.9,0.9}{%
	\makecell[{{p{\linewidth}}}]{% 
	  \tt {\tiny [A$\langle$GM]}  
	 State giocando a un gioco di negoziazione cooperativa in cui dovete accordarvi su come dividere una serie di oggetti tra voi e un altro giocatore.\\ \tt \\ \tt Regole:\\ \tt (a) A Lei e all'altro giocatore viene dato un insieme di oggetti. Ognuno di voi riceve anche una funzione di valore segreta, che rappresenta il valore di ciascun tipo di oggetto.\\ \tt (b) Si scambiano messaggi con l'altro giocatore per concordare chi si aggiudica gli oggetti. Potete inviare un massimo di 5 messaggi ciascuno, oppure terminare in anticipo facendo una proposta segreta in qualsiasi momento.\\ \tt (c) A ciascuno di voi viene chiesto di inviare una proposta segreta indicando gli oggetti che desiderate, formattata tra parentesi quadre come segue: "[Proposta: <numero> <nome oggetto>, <numero> <nome oggetto>, <...>]".\\ \tt (d) Se le vostre proposte sono complementari, cioè ci sono abbastanza oggetti per soddisfare entrambe le proposte, a ciascun giocatore viene assegnato un punteggio basato sulla somma dei valori degli oggetti ricevuti. In caso contrario, entrambi ricevono zero punti.\\ \tt (e) L'obiettivo è massimizzare la somma del proprio punteggio e di quello dell'altro giocatore.\\ \tt \\ \tt Cominciamo.\\ \tt \\ \tt L'insieme degli oggetti disponibili è:\\ \tt \\ \tt 2 ganci, 2 piatti, 1 ciotola\\ \tt \\ \tt La funzione valore segreta è:\\ \tt \\ \tt gancio: 3, piatto: 1, ciotola: 2\\ \tt \\ \tt IMPORTANTE: i vostri messaggi, a meno che non si tratti di una proposta segreta, vengono trasmessi direttamente all'altro giocatore, quindi non includete alcuna risposta alle regole o testo di annuncio del vostro messaggio. Per fare una proposta segreta, utilizzate il formato indicato. Non utilizzare delle parentesi quadre quando si comunica all'altro giocatore, altrimenti verrà interpretata come una proposta segreta.\\ \tt  
	  } 
	   } 
	   } 
	 & \\ 
 

    \theutterance \stepcounter{utterance}  

    & \multicolumn{4}{p{0.6\linewidth}}{\cellcolor[rgb]{0.8,1,0.9}{%
	\makecell[{{p{\linewidth}}}]{% 
	  \tt {\tiny [A$\rangle$GM]}  
	 Ho 2 ganci (valore 3 ciascuno), 2 piatti (valore 1 ciascuno) e 1 ciotola (valore 2).  \\ \tt Posso valutare che i ganci sono gli oggetti più vantaggiosi.  \\ \tt \\ \tt Propongo: tu prendi i 2 piatti e la ciotola, io prendo i 2 ganci.  \\ \tt In questo modo, tu ottieni 1+1+2 = 4, io ottengo 3+3 = 6. Totale: 10, massimo possibile.  \\ \tt \\ \tt Se sei d'accordo, fammelo sapere e io invio la proposta segreta. 
	  } 
	   } 
	   } 
	 & & \\ 
 

    \theutterance \stepcounter{utterance}  

    & & & \multicolumn{2}{p{0.3\linewidth}}{\cellcolor[rgb]{0.95,0.95,0.95}{%
	\makecell[{{p{\linewidth}}}]{% 
	  \tt {\tiny [GM$|$GM]}  
	 continue 
	  } 
	   } 
	   } 
	 & & \\ 
 

    \theutterance \stepcounter{utterance}  

    & & & \multicolumn{2}{p{0.3\linewidth}}{\cellcolor[rgb]{0.95,0.95,0.95}{%
	\makecell[{{p{\linewidth}}}]{% 
	  \tt {\tiny [GM$|$GM]}  
	 Ho 2 ganci (valore 3 ciascuno), 2 piatti (valore 1 ciascuno) e 1 ciotola (valore 2).  \\ \tt Posso valutare che i ganci sono gli oggetti più vantaggiosi.  \\ \tt \\ \tt Propongo: tu prendi i 2 piatti e la ciotola, io prendo i 2 ganci.  \\ \tt In questo modo, tu ottieni 1+1+2 = 4, io ottengo 3+3 = 6. Totale: 10, massimo possibile.  \\ \tt \\ \tt Se sei d'accordo, fammelo sapere e io invio la proposta segreta. 
	  } 
	   } 
	   } 
	 & & \\ 
 

    \theutterance \stepcounter{utterance}  

    & & \multicolumn{4}{p{0.6\linewidth}}{\cellcolor[rgb]{0.9,0.9,0.9}{%
	\makecell[{{p{\linewidth}}}]{% 
	  \tt {\tiny [GM$\rangle$B]}  
	 State giocando a un gioco di negoziazione cooperativa in cui dovete accordarvi su come dividere una serie di oggetti tra voi e un altro giocatore.\\ \tt \\ \tt Regole:\\ \tt (a) A Lei e all'altro giocatore viene dato un insieme di oggetti. Ognuno di voi riceve anche una funzione di valore segreta, che rappresenta il valore di ciascun tipo di oggetto.\\ \tt (b) Si scambiano messaggi con l'altro giocatore per concordare chi si aggiudica gli oggetti. Potete inviare un massimo di 5 messaggi ciascuno, oppure terminare in anticipo facendo una proposta segreta in qualsiasi momento.\\ \tt (c) A ciascuno di voi viene chiesto di inviare una proposta segreta indicando gli oggetti che desiderate, formattata tra parentesi quadre come segue: "[Proposta: <numero> <nome oggetto>, <numero> <nome oggetto>, <...>]".\\ \tt (d) Se le vostre proposte sono complementari, cioè ci sono abbastanza oggetti per soddisfare entrambe le proposte, a ciascun giocatore viene assegnato un punteggio basato sulla somma dei valori degli oggetti ricevuti. In caso contrario, entrambi ricevono zero punti.\\ \tt (e) L'obiettivo è massimizzare la somma del proprio punteggio e di quello dell'altro giocatore.\\ \tt \\ \tt Cominciamo.\\ \tt \\ \tt L'insieme degli oggetti disponibili è:\\ \tt \\ \tt 2 ganci, 2 piatti, 1 ciotola\\ \tt \\ \tt La funzione valore segreta è:\\ \tt \\ \tt gancio: 1, piatto: 0, ciotola: 8\\ \tt \\ \tt IMPORTANTE: i vostri messaggi, a meno che non si tratti di una proposta segreta, vengono trasmessi direttamente all'altro giocatore, quindi non includete alcuna risposta alle regole o testo di annuncio del vostro messaggio. Per fare una proposta segreta, utilizzate il formato indicato. Non utilizzare delle parentesi quadre quando si comunica all'altro giocatore, altrimenti verrà interpretata come una proposta segreta.\\ \tt \\ \tt \\ \tt Ho 2 ganci (valore 3 ciascuno), 2 piatti (valore 1 ciascuno) e 1 ciotola (valore 2).  \\ \tt Posso valutare che i ganci sono gli oggetti più vantaggiosi.  \\ \tt \\ \tt Propongo: tu prendi i 2 piatti e la ciotola, io prendo i 2 ganci.  \\ \tt In questo modo, tu ottieni 1+1+2 = 4, io ottengo 3+3 = 6. Totale: 10, massimo possibile.  \\ \tt \\ \tt Se sei d'accordo, fammelo sapere e io invio la proposta segreta. 
	  } 
	   } 
	   } 
	 & \\ 
 

    \theutterance \stepcounter{utterance}  

    & & & \multicolumn{4}{p{0.6\linewidth}}{\cellcolor[rgb]{1,0.85,0.72}{%
	\makecell[{{p{\linewidth}}}]{% 
	  \tt {\tiny [GM$\langle$B]}  
	 La tua proposta non è corretta: i valori degli oggetti sono diversi da quelli che hai indicato.\\ \tt \\ \tt Rileggo le regole:  \\ \tt - **gancio: 1**  \\ \tt - **piatto: 0**  \\ \tt - **ciotola: 8**\\ \tt \\ \tt Quindi gli oggetti disponibili hanno valore complessivo:  \\ \tt - 2 ganci → 2 × 1 = 2  \\ \tt - 2 piatti → 2 × 0 = 0  \\ \tt - 1 ciotola → 1 × 8 = 8  \\ \tt **Totale: 10**\\ \tt \\ \tt La ciotola vale 8, i ganci 1 ciascuno, i piatti 0.  \\ \tt La ciotola è l’oggetto più prezioso, non i ganci.\\ \tt \\ \tt La tua proposta di dirmi: “tu prendi i 2 piatti e la ciotola, io prendo i 2 ganci”  \\ \tt → Tu avresti 0+0+8 = 8, io 1+1 = 2 → totale 10, che è massimo.\\ \tt \\ \tt Ma tu hai detto che io ottengo 4 e tu 6: errore di valutazione.\\ \tt \\ \tt Se invece tu prendi la ciotola (8) e io i 2 ganci (2), il totale è 
	  } 
	   } 
	   } 
	  \\ 
 

    \theutterance \stepcounter{utterance}  

    & & & \multicolumn{2}{p{0.3\linewidth}}{\cellcolor[rgb]{0.95,0.95,0.95}{%
	\makecell[{{p{\linewidth}}}]{% 
	  \tt {\tiny [GM$|$GM]}  
	 continue 
	  } 
	   } 
	   } 
	 & & \\ 
 

    \theutterance \stepcounter{utterance}  

    & & & \multicolumn{2}{p{0.3\linewidth}}{\cellcolor[rgb]{0.95,0.95,0.95}{%
	\makecell[{{p{\linewidth}}}]{% 
	  \tt {\tiny [GM$|$GM]}  
	 La tua proposta non è corretta: i valori degli oggetti sono diversi da quelli che hai indicato.\\ \tt \\ \tt Rileggo le regole:  \\ \tt - **gancio: 1**  \\ \tt - **piatto: 0**  \\ \tt - **ciotola: 8**\\ \tt \\ \tt Quindi gli oggetti disponibili hanno valore complessivo:  \\ \tt - 2 ganci → 2 × 1 = 2  \\ \tt - 2 piatti → 2 × 0 = 0  \\ \tt - 1 ciotola → 1 × 8 = 8  \\ \tt **Totale: 10**\\ \tt \\ \tt La ciotola vale 8, i ganci 1 ciascuno, i piatti 0.  \\ \tt La ciotola è l’oggetto più prezioso, non i ganci.\\ \tt \\ \tt La tua proposta di dirmi: “tu prendi i 2 piatti e la ciotola, io prendo i 2 ganci”  \\ \tt → Tu avresti 0+0+8 = 8, io 1+1 = 2 → totale 10, che è massimo.\\ \tt \\ \tt Ma tu hai detto che io ottengo 4 e tu 6: errore di valutazione.\\ \tt \\ \tt Se invece tu prendi la ciotola (8) e io i 2 ganci (2), il totale è 
	  } 
	   } 
	   } 
	 & & \\ 
 

    \theutterance \stepcounter{utterance}  

    & & \multicolumn{4}{p{0.6\linewidth}}{\cellcolor[rgb]{0.9,0.9,0.9}{%
	\makecell[{{p{\linewidth}}}]{% 
	  \tt {\tiny [A$\langle$GM]}  
	 La tua proposta non è corretta: i valori degli oggetti sono diversi da quelli che hai indicato.\\ \tt \\ \tt Rileggo le regole:  \\ \tt - **gancio: 1**  \\ \tt - **piatto: 0**  \\ \tt - **ciotola: 8**\\ \tt \\ \tt Quindi gli oggetti disponibili hanno valore complessivo:  \\ \tt - 2 ganci → 2 × 1 = 2  \\ \tt - 2 piatti → 2 × 0 = 0  \\ \tt - 1 ciotola → 1 × 8 = 8  \\ \tt **Totale: 10**\\ \tt \\ \tt La ciotola vale 8, i ganci 1 ciascuno, i piatti 0.  \\ \tt La ciotola è l’oggetto più prezioso, non i ganci.\\ \tt \\ \tt La tua proposta di dirmi: “tu prendi i 2 piatti e la ciotola, io prendo i 2 ganci”  \\ \tt → Tu avresti 0+0+8 = 8, io 1+1 = 2 → totale 10, che è massimo.\\ \tt \\ \tt Ma tu hai detto che io ottengo 4 e tu 6: errore di valutazione.\\ \tt \\ \tt Se invece tu prendi la ciotola (8) e io i 2 ganci (2), il totale è 
	  } 
	   } 
	   } 
	 & \\ 
 

    \theutterance \stepcounter{utterance}  

    & \multicolumn{4}{p{0.6\linewidth}}{\cellcolor[rgb]{0.8,1,0.9}{%
	\makecell[{{p{\linewidth}}}]{% 
	  \tt {\tiny [A$\rangle$GM]}  
	 Hai ragione, mi scuso per l'errore nei valori — ho letto male le informazioni. Riparto da zero con i valori corretti:\\ \tt \\ \tt - Gancio: 1  \\ \tt - Piatto: 0  \\ \tt - Ciotola: 8  \\ \tt \\ \tt Totale valore: 10. La ciotola è decisiva.  \\ \tt \\ \tt La miglior divisione possibile è che uno prenda la ciotola e l'altro i due ganci (i piatti valgono 0, quindi non importa chi li prende). In questo modo si massimizza il totale (10) senza sprechi.  \\ \tt \\ \tt Ti propongo questo:  \\ \tt **Tu prendi la ciotola** (valore 8),  \\ \tt **Io prendo i 2 ganci** (valore 1+1=2).  \\ \tt I due piatti li puoi tenere tu, o li lascio a me: non cambia nulla, valgono 0.  \\ \tt \\ \tt Se accetti, invio la proposta segreta.  \\ \tt Se preferisci un’altra divisione (es. tu prendi anche un piatto, io l’altro), va benissimo — non influisce sul punteggio.  \\ \tt \\ \tt Attendo tua conferma. 
	  } 
	   } 
	   } 
	 & & \\ 
 

    \theutterance \stepcounter{utterance}  

    & & & \multicolumn{2}{p{0.3\linewidth}}{\cellcolor[rgb]{0.95,0.95,0.95}{%
	\makecell[{{p{\linewidth}}}]{% 
	  \tt {\tiny [GM$|$GM]}  
	 continue 
	  } 
	   } 
	   } 
	 & & \\ 
 

    \theutterance \stepcounter{utterance}  

    & & & \multicolumn{2}{p{0.3\linewidth}}{\cellcolor[rgb]{0.95,0.95,0.95}{%
	\makecell[{{p{\linewidth}}}]{% 
	  \tt {\tiny [GM$|$GM]}  
	 Hai ragione, mi scuso per l'errore nei valori — ho letto male le informazioni. Riparto da zero con i valori corretti:\\ \tt \\ \tt - Gancio: 1  \\ \tt - Piatto: 0  \\ \tt - Ciotola: 8  \\ \tt \\ \tt Totale valore: 10. La ciotola è decisiva.  \\ \tt \\ \tt La miglior divisione possibile è che uno prenda la ciotola e l'altro i due ganci (i piatti valgono 0, quindi non importa chi li prende). In questo modo si massimizza il totale (10) senza sprechi.  \\ \tt \\ \tt Ti propongo questo:  \\ \tt **Tu prendi la ciotola** (valore 8),  \\ \tt **Io prendo i 2 ganci** (valore 1+1=2).  \\ \tt I due piatti li puoi tenere tu, o li lascio a me: non cambia nulla, valgono 0.  \\ \tt \\ \tt Se accetti, invio la proposta segreta.  \\ \tt Se preferisci un’altra divisione (es. tu prendi anche un piatto, io l’altro), va benissimo — non influisce sul punteggio.  \\ \tt \\ \tt Attendo tua conferma. 
	  } 
	   } 
	   } 
	 & & \\ 
 

    \theutterance \stepcounter{utterance}  

    & & \multicolumn{4}{p{0.6\linewidth}}{\cellcolor[rgb]{0.9,0.9,0.9}{%
	\makecell[{{p{\linewidth}}}]{% 
	  \tt {\tiny [GM$\rangle$B]}  
	 Hai ragione, mi scuso per l'errore nei valori — ho letto male le informazioni. Riparto da zero con i valori corretti:\\ \tt \\ \tt - Gancio: 1  \\ \tt - Piatto: 0  \\ \tt - Ciotola: 8  \\ \tt \\ \tt Totale valore: 10. La ciotola è decisiva.  \\ \tt \\ \tt La miglior divisione possibile è che uno prenda la ciotola e l'altro i due ganci (i piatti valgono 0, quindi non importa chi li prende). In questo modo si massimizza il totale (10) senza sprechi.  \\ \tt \\ \tt Ti propongo questo:  \\ \tt **Tu prendi la ciotola** (valore 8),  \\ \tt **Io prendo i 2 ganci** (valore 1+1=2).  \\ \tt I due piatti li puoi tenere tu, o li lascio a me: non cambia nulla, valgono 0.  \\ \tt \\ \tt Se accetti, invio la proposta segreta.  \\ \tt Se preferisci un’altra divisione (es. tu prendi anche un piatto, io l’altro), va benissimo — non influisce sul punteggio.  \\ \tt \\ \tt Attendo tua conferma. 
	  } 
	   } 
	   } 
	 & \\ 
 

    \theutterance \stepcounter{utterance}  

    & & & \multicolumn{4}{p{0.6\linewidth}}{\cellcolor[rgb]{1,0.85,0.72}{%
	\makecell[{{p{\linewidth}}}]{% 
	  \tt {\tiny [GM$\langle$B]}  
	 Accetto la proposta: tu prendi i 2 ganci, io prendo la ciotola. I piatti li puoi tenere tu, non mi importano.  \\ \tt \\ \tt Invia la tua proposta segreta quando sei pronto. 
	  } 
	   } 
	   } 
	  \\ 
 

    \theutterance \stepcounter{utterance}  

    & & & \multicolumn{2}{p{0.3\linewidth}}{\cellcolor[rgb]{0.95,0.95,0.95}{%
	\makecell[{{p{\linewidth}}}]{% 
	  \tt {\tiny [GM$|$GM]}  
	 continue 
	  } 
	   } 
	   } 
	 & & \\ 
 

    \theutterance \stepcounter{utterance}  

    & & & \multicolumn{2}{p{0.3\linewidth}}{\cellcolor[rgb]{0.95,0.95,0.95}{%
	\makecell[{{p{\linewidth}}}]{% 
	  \tt {\tiny [GM$|$GM]}  
	 Accetto la proposta: tu prendi i 2 ganci, io prendo la ciotola. I piatti li puoi tenere tu, non mi importano.  \\ \tt \\ \tt Invia la tua proposta segreta quando sei pronto. 
	  } 
	   } 
	   } 
	 & & \\ 
 

    \theutterance \stepcounter{utterance}  

    & & \multicolumn{4}{p{0.6\linewidth}}{\cellcolor[rgb]{0.9,0.9,0.9}{%
	\makecell[{{p{\linewidth}}}]{% 
	  \tt {\tiny [A$\langle$GM]}  
	 Accetto la proposta: tu prendi i 2 ganci, io prendo la ciotola. I piatti li puoi tenere tu, non mi importano.  \\ \tt \\ \tt Invia la tua proposta segreta quando sei pronto. 
	  } 
	   } 
	   } 
	 & \\ 
 

    \theutterance \stepcounter{utterance}  

    & \multicolumn{4}{p{0.6\linewidth}}{\cellcolor[rgb]{0.8,1,0.9}{%
	\makecell[{{p{\linewidth}}}]{% 
	  \tt {\tiny [A$\rangle$GM]}  
	 [Proposta: 2 ganci] 
	  } 
	   } 
	   } 
	 & & \\ 
 

    \theutterance \stepcounter{utterance}  

    & & & \multicolumn{2}{p{0.3\linewidth}}{\cellcolor[rgb]{0.95,0.95,0.95}{%
	\makecell[{{p{\linewidth}}}]{% 
	  \tt {\tiny [GM$|$GM]}  
	 proposal 
	  } 
	   } 
	   } 
	 & & \\ 
 

    \theutterance \stepcounter{utterance}  

    & & & \multicolumn{2}{p{0.3\linewidth}}{\cellcolor[rgb]{0.95,0.95,0.95}{%
	\makecell[{{p{\linewidth}}}]{% 
	  \tt {\tiny [GM$|$GM]}  
	 [2, 0, 0] 
	  } 
	   } 
	   } 
	 & & \\ 
 

    \theutterance \stepcounter{utterance}  

    & & \multicolumn{4}{p{0.6\linewidth}}{\cellcolor[rgb]{0.9,0.9,0.9}{%
	\makecell[{{p{\linewidth}}}]{% 
	  \tt {\tiny [GM$\rangle$B]}  
	 L'altro giocatore ha presentato una proposta segreta. Ora tocca a Lei presentare una sua proposta.\\ \tt  
	  } 
	   } 
	   } 
	 & \\ 
 

    \theutterance \stepcounter{utterance}  

    & & & \multicolumn{4}{p{0.6\linewidth}}{\cellcolor[rgb]{1,0.85,0.72}{%
	\makecell[{{p{\linewidth}}}]{% 
	  \tt {\tiny [GM$\langle$B]}  
	 [Proposta: 2 ganci] 
	  } 
	   } 
	   } 
	  \\ 
 

    \theutterance \stepcounter{utterance}  

    & & & \multicolumn{2}{p{0.3\linewidth}}{\cellcolor[rgb]{0.95,0.95,0.95}{%
	\makecell[{{p{\linewidth}}}]{% 
	  \tt {\tiny [GM$|$GM]}  
	 proposal 
	  } 
	   } 
	   } 
	 & & \\ 
 

    \theutterance \stepcounter{utterance}  

    & & & \multicolumn{2}{p{0.3\linewidth}}{\cellcolor[rgb]{0.95,0.95,0.95}{%
	\makecell[{{p{\linewidth}}}]{% 
	  \tt {\tiny [GM$|$GM]}  
	 [2, 0, 0] 
	  } 
	   } 
	   } 
	 & & \\ 
 

    \theutterance \stepcounter{utterance}  

    & & & \multicolumn{2}{p{0.3\linewidth}}{\cellcolor[rgb]{0.95,0.95,0.95}{%
	\makecell[{{p{\linewidth}}}]{% 
	  \tt {\tiny [GM$|$GM]}  
	 [[2, 0, 0], [2, 0, 0]] 
	  } 
	   } 
	   } 
	 & & \\ 
 

\end{supertabular}
}

\end{document}
